\chapter{Sprint 4: Automation, Evaluation, and Optimization}
\chaptertoc

\clearpage

\section{Introduction}
Sprint 4 completed the monitoring loop by connecting queue alerts to an external notification channel and validating the system as a whole. Until this sprint, the project could detect queues, estimate waiting time, expose live services, and display the dashboard. The remaining step was to make alerts reach operational staff and to verify that the full chain behaved correctly under realistic conditions.

The sprint therefore combined two types of work: alert automation through n8n and Telegram, and final evaluation of perception, live monitoring, transport fallback, and notification delivery.

\section{Sprint Objectives and Backlog}
Sprint 4 focused on alert automation and system-wide verification. Alert automation required a validated payload, secure workflow access, cooldown handling, and Telegram delivery. Verification required testing the integrated system, checking streaming stability, tuning runtime responsiveness, and improving recovery from temporary failures. The two main manager needs were alert automation and reliable notification delivery. Table~\ref{tab:sprint4_user_stories} lists the user stories that drove this work.

\begin{table}[H]
    \centering
    \begingroup\setlength{\tabcolsep}{4pt}\renewcommand{\arraystretch}{0.95}\footnotesize
    \begin{tabularx}{\textwidth}{|l|X|l|X|}
        \hline
        User Story ID & User Story & Task ID & Task \\ \hline
        US4.1 & As a manager, I want alert payloads dispatched to n8n so that operational staff can be notified automatically. & T4.1 & Define the webhook payload schema and validate it before dispatch. \\ \hline
        US4.3 & As a manager, I want alerts routed through Telegram so that the on-site operator receives notifications on a familiar channel. & T4.3 & Create and configure the Telegram bot, store its credential in n8n, and configure the workflow. \\ \hline
    \end{tabularx}
    \caption{Sprint 4 user stories}
    \label{tab:sprint4_user_stories}
    \endgroup
\end{table}

\subsection{Technical Tasks}
Beyond the user stories, the sprint included reliability, testing, and performance work. The technical breakdown is shown in Table~\ref{tab:sprint4_technical_tasks}.

\begin{table}[H]
    \centering
    \begingroup\setlength{\tabcolsep}{4pt}\renewcommand{\arraystretch}{0.95}\footnotesize
    \begin{tabularx}{\textwidth}{|l|X|l|X|}
        \hline
        Task ID & Task & Subtask ID & Subtask \\ \hline
        T4.2 & Add retry logic, timeout handling, and logging for webhook delivery to ensure notifications are not silently lost. & ST4.2.1 & Implement delivery resilience with exponential backoff and failure tracking. \\ \hline
        T4.4 & Prepare PowerShell scripts for single-alert, multi-feed, and execution polling scenarios to validate n8n workflow behavior. & ST4.4.1 & Create repeatable test automation for workflow routing and alerting paths. \\ \hline
        T4.5 & Run backend and frontend tests together and confirm lifecycle, streaming, and API contracts under realistic conditions. & ST4.5.1 & Execute comprehensive integration verification to catch issues before delivery. \\ \hline
        T4.6 & Apply runtime tuning: configurable image size, frame stride, async dispatch, and playback pacing to improve processing responsiveness. & ST4.6.1 & Implement performance optimizations for the main processing loop. \\ \hline
        T4.7 & Improve feed worker supervision, frame handoff continuity, and websocket reconnection to strengthen fault tolerance. & ST4.7.1 & Enhance system recovery from transient failures and network disruptions. \\ \hline
    \end{tabularx}
    \caption{Sprint 4 technical tasks}
    \label{tab:sprint4_technical_tasks}
    \endgroup
\end{table}

\section{Design}

\subsection{Alert Automation and n8n Workflow}
The alert workflow was designed to avoid two opposite problems: missing important queue alerts and overwhelming staff with repeated notifications for the same condition. When an alert arrives, the workflow first checks that the message is authorized. It then applies cooldown logic, archives the alert in the database, and sends the validated notification to Telegram.

This separation makes the workflow easier to reason about. Authorization protects the entry point, cooldown controls repetition, archival supports traceability, and Telegram delivery reaches the on-site staff channel. Figure~\ref{fig:sprint4_workflow} illustrates the complete workflow logic.

\begin{figure}[!htbp]
    \centering
    \includegraphics[width=\textwidth]{\detokenize{n8n_workflow.png}}
    \caption{Sprint 4 n8n workflow for routing queue alerts through cooldown checks, archival, and Telegram delivery.}
    \label{fig:sprint4_workflow}
\end{figure}

\subsection{System Verification and Evaluation Strategy}
The evaluation strategy covered the main parts of the finished system: queue analytics, API behavior, dashboard responsiveness, transport fallback, and n8n routing. The most important observations were whether the perception pipeline stayed responsive, whether streams recovered after feed changes, whether alerts reached Telegram, and whether the dashboard stayed usable when WebSocket delivery was interrupted.

The tests combined focused checks for individual behaviors with integrated scenarios that followed the system from queue-state change to dashboard update or Telegram notification. Normal operation, temporary failures, feed lifecycle changes, and alert cooldown behavior were all included because these are the situations most likely to affect a real monitoring deployment.

Figure~\ref{fig:sprint4_evaluation_resilience_activity} illustrates the iterative evaluation and tuning cycle that characterized the sprint.

\begin{figure}[!htbp]
    \centering
    \includegraphics[width=0.74\textwidth]{\detokenize{chapter6_evaluation_resilience_activity.pdf}}
    \caption{Sprint 4 evaluation and resilience activity diagram for tests, tuning, and recovery verification.}
    \label{fig:sprint4_evaluation_resilience_activity}
\end{figure}

\section{Implementation}

\subsection{Alert Payload and Delivery Architecture}
The alert payload contains the information needed by the automation workflow and the Telegram message: timestamp and frame context, queue metrics ($\lambda$, $\mu$, and expected wait time), alert flags, and severity. Payload validation happens before dispatch so incomplete or malformed messages are not sent to external workflow tools.

Delivery reliability was handled through retry behavior, timeouts, shared-secret validation, and dispatch logging. Successful and failed delivery attempts are kept for traceability, which is important when a manager needs to know whether a threshold violation was actually sent to the notification workflow.

\subsection{Telegram Integration and Notification Delivery}
A Telegram bot was configured as the notification channel, with its credential stored in n8n rather than directly inside the exported workflow. This keeps the workflow portable and avoids exposing secrets in shared project artifacts. Alerts that pass authorization and cooldown checks are formatted for human readability and sent to the operational group chat. An example of delivered alert formatting is shown in Figure~\ref{fig:sprint4_telegram_alert}.

\begin{figure}[!htbp]
    \centering
    \includegraphics[width=0.62\textwidth]{\detokenize{telegram_alert_screenshot.png}}
    \caption{Example Telegram alert delivered.}
    \label{fig:sprint4_telegram_alert}
\end{figure}

\subsection{Asynchronous Alert Dispatch and Deduplication}
Alert delivery was separated from the perception loop so that network latency would not slow frame processing. This matters because webhook calls depend on external services and can be delayed or fail temporarily. Queue analysis should continue even if the notification path is slower than expected.

Deduplication reduces repeated notifications when the same threshold violation remains active across several metric windows. A new alert is still sent when the condition changes or a new threshold violation appears, but persistent conditions do not flood the Telegram group.

\subsection{n8n Workflow Structure and Validation}
The n8n workflow separates three alert states: no alert, suppressed alert during cooldown, and delivery-ready alert. This structure made validation clearer because each branch could be exercised independently.

PowerShell validation scripts were used for single-alert scenarios, multi-feed scenarios, and execution polling. These checks confirmed that authorization, cooldown routing, and Telegram delivery behaved as expected across different alert paths.

\subsection{Comprehensive Integration Verification}
Integration verification covered the complete monitoring chain: feed lifecycle actions, streaming behavior, live dashboard updates, webhook delivery, and Telegram notification routing. Error scenarios included authentication failures, WebSocket disconnections, WebRTC unavailability, and feed worker interruption.

The goal was not only to confirm that each feature worked in isolation, but also to check that failures in one part of the system did not make the whole platform unusable. For example, a WebSocket interruption should move the dashboard to polling, and a temporary notification failure should be recorded rather than silently ignored.

\subsection{Runtime Optimization and Performance Tuning}
Several runtime tuning options were added to improve responsiveness. Image size and frame stride can be adjusted so the project can trade some detection detail for faster processing when hardware resources are limited. Device selection supports different deployment machines. Alert dispatch was moved away from the main perception path, and file playback pacing was adjusted to better simulate real-time delivery instead of pushing frames as fast as possible.

\subsection{Fault-Tolerance and Recovery Mechanisms}
Fault tolerance was improved at several points in the system. Feed supervision helps recover from worker interruption, frame handoff was hardened to reduce loss during transitions, and the dashboard reconnects automatically before falling back to polling. Together, these behaviors keep monitoring available during temporary failures rather than requiring a full manual restart for every interruption.

\subsection{System Integration and Operational Readiness}
By the end of Sprint 4, the perception pipeline, queue analytics, backend services, web dashboard, and n8n alert automation were connected into one working platform. The remaining work before production deployment is mainly deeper long-duration testing, load evaluation, and infrastructure hardening, not the completion of missing core features.

\section{Conclusion}
Sprint 4 delivered the final monitoring-to-action loop. Queue alerts could be validated, routed through n8n, filtered by cooldown logic, and delivered to Telegram. The evaluation work confirmed the main integrated behaviors: live metrics, dashboard resilience, stream fallback, alert dispatch, and workflow routing.

The sprint also showed why end-to-end validation is necessary for this type of project. Component checks can confirm isolated behavior, but only integrated scenarios reveal whether the manager still receives useful information when transport, notification, or feed conditions degrade. After this sprint, the project reached a functional level suitable for demonstration and further production-oriented validation.

